\documentclass{article}
\usepackage{geometry}
\usepackage{concmath}
\usepackage[OT1]{fontenc}
\usepackage{graphicx}
\usepackage{amsfonts}
\usepackage{amsmath,mathtools}
\usepackage{calrsfs}
\usepackage{color}
\usepackage[dvipsnames]{xcolor}
\usepackage[many]{tcolorbox}
\usepackage{setspace}
\usepackage{tkz-tab}    
\usepackage{amssymb}          
\usepackage{multicol}  
\usepackage{esint}
\geometry{margin=0.5in}
\definecolor{green}{RGB}{52, 235, 168}
\definecolor{violet}{RGB}{165, 83, 237}
\definecolor{lemon}{RGB}{170, 237, 83}
\definecolor{red}{RGB}{240, 125, 84}
\definecolor{blue}{RGB}{105, 255, 242}

\begin{document} \noindent
\colorbox{green}{260.1}\\
Soit \(M \in \mathcal{M}_n(\mathbb{R})\) tel que \begin{align*}
    &M^2 - M^T = I_n \\
    &(M^T)^2 - M = I_n \\
    &(M^2 - I_n)^2 - M = I_n \\
    &M^4 - 2M^2 + I_n - M = I_n \\
    P &= X^4 - 2X^2 - X \text{ annule M}\\
    &= X(X^3-2-1)\\
    &= X(X+1)(X^2 - X -1)
\end{align*}
Donc P est scindé à racines simples \\
Donc \fbox{M est diagonalisable}\\

\noindent \colorbox{green}{260.16}\\
\begin{align*}
    M &= \begin{pmatrix}
        0&-1& & (0) & \\
        1 & 0 & & \\
         & & 0 & 0 & 1\\
        &(0)& 1 & 0 & 0 \\
        & & 0 & 1 & 0
    \end{pmatrix}\\
    M^k &= \begin{pmatrix}
        A^k & 0 \\
        0 & B^k
    \end{pmatrix}\\
    A^2 &= \begin{pmatrix}
        -1 & 0 \\ 0 & -1
    \end{pmatrix}\\
    A^4 &= I_2\\
    B^2 &= \begin{pmatrix}
        0 & 1 & 0 \\ 0 & 0 & 1\\ 1& 0 & 0
    \end{pmatrix}\\
    B^3 &= I_3\\
    p &= ppcm(4,3) = 12 \text{ constant }\\
    Q &= X^12-1 \text{ est annulateur de M}\\
    rac(Q) = \mathbb{U}_{12}
\end{align*}
Q scindé à racines simples. donc M est diagonalisable\\
b/ \begin{align*}
    x = e_1 + e_3\\
    m(x) = e_2 + e_4 \\
    m^2 (x) = e_1 + e_5\\
    m^3(x) = -e_2 + e_3\\
    m^4(x) = e_1 + e_4 \\
    m^5(x) = e_2 + e_5 
\end{align*}
B base canonique : \begin{align*}
    det(x, m(x),...,m^4(x)) &= \begin{vmatrix}
        1 & 0 & -1 & 0 & 1\\
        0 & 1 & 0 & -1 & 0\\
        1& 0& 0 & 1 & 0 \\
        0& 1& 0 & 0 & 1\\
        0& 0& 1& 0 & 0\\
    \end{vmatrix}\\
    &= 1 \begin{vmatrix}
        1 & 0 & 0 & 1\\
        0 & 1 & -1 & 0\\
        1& 0 & 1 & 0 \\
        0& 1 & 0 & -1
    \end{vmatrix}\\
    &=1\begin{vmatrix}
        1 & 0 & 1 \\ 0 & -1& 0 \\ 1&1&0
    \end{vmatrix}
    + 1 \begin{vmatrix}
        1 & 0 & 0 \\ 0 & 1 & -1 \\ 1& 0 & 1
    \end{vmatrix}\\
    &= 1\times 1 + 1 \times 1\\
    &= 2 \\
    &\neq 0\\
    M' &= \begin{pmatrix}
        0 & 0 & 0 & 0 & 1\\
        1 & 0 & 0 & 0 & 0 \\
        0 & 1 & 0 & 0 & 1\\
        0 & 0 & 1 & 0 & -1\\
        0 & 0 & 0 & 1 & 0
    \end{pmatrix}
\end{align*}

\noindent \colorbox{green}{260.13}\\
rgA = 1, donc ImA est une droite vectorielle \begin{equation*}
    \exists X \in M_{n,1}(\mathbb{R}), \text{ tq } ImA = Vect(X)
\end{equation*}
* \(C_k \in ImA\) donc \(\exists y_k\) tq  \(C_k = y_k X\)\\
Posons \(Y = \begin{pmatrix}
    y_1\\\vdots\\y_n
\end{pmatrix}\) on a \(A = XY^T\)
\begin{align*}
    A^2 &= XY^TXY^T\\
    &= X(Y^TX)^TY^T\\
    &= \lambda XY^T\\
    &= \lambda A
\end{align*}
Avec \(\lambda = \sum_{k=1}^{n}x_ky_k \Leftrightarrow \lambda = trA\)\\
b/ Si A diagonalisable. \(\exists Q \in GL_n(\mathbb{K})\) 
tq A = \(Q\begin{pmatrix}
    0 & & (0)\\
    & \ddots & \\
    (0) & & 0
\end{pmatrix}Q^{-1}\)\\
Or rg A = 1\\
On a \(X^2\) annulateur de A\\
\(\Pi_A | X^2\) donc \(\Pi_A = X\) ou \(\Pi_A = X^2\)\\
\underline{Si \(\Pi_A = X\)} X annule A, donc A  = 0. Absurde\\
Donc \(\Pi_A = X^2\), or \(\Pi_A\) n'est pas scindé à racines simples
donc A n'est pas diagonalisable\\

\noindent \colorbox{green}{260.4}\\
\(u \in \mathcal{L}(E)\) tq \(u^3 = Id\\
P = X^3 - 1\)\\
P annulateur de u \begin{align*}
    P &= (X-1)(X^2 + X + 1)\\
    (X-1)&\wedge(X^2 + X + 1) =1\\
    Ker(P(u)) &= Ker(P_1(u))\oplus Ker(P_2(u))\\
    E &= Ker(u-Id) \oplus Ker(u^2 + u + Id)\\
    Im(u-Id) &= Ker(u^2 + u + Id)\\
    rg(P_1(u)) &= dimE - dimKer(P_1(u))\\
    &= dimKer(u^2 + u + id)
\end{align*}
Soit \(x \in Im(P_1(u))\)\begin{align*}
    \exists t \in E \text{ tq } x &= P_1(u)(t)\\
    P_2(u)(x) &= P_2(u) \circ P_1(u)(t)\\
    &= (P_2P_1)(u)(t)\\
    &= P(u)(t)\\
    &= 0 
\end{align*}

\noindent \colorbox{green}{260.3}\\
\begin{align*}
    A = \begin{pmatrix}
        1 & & ... & & 1\\
         & 0 & ... & 0 & \\
        \vdots & \vdots &  & \vdots & \vdots \\
         & 0 & ... & 0 & \\
        1 & & ... & & 1
    \end{pmatrix}
\end{align*}
1/ rgA = 2\\
2/ dimKerA = n-2\\
\(dimE_0(A) = n-2 \leq m_0 \)\\
\(dimE_0(A) = n-2 = m_0 \) car A symétrique\\
3/ \begin{align*}
    A^2 &= \begin{pmatrix}
        n &2 & ... &2& n\\
         2& 2 & ... & 2 & 2\\
        \vdots & \vdots &  & \vdots & \vdots \\
         2& 2 & ... & 2 & 2\\
        n& 2& ... &2 & n
    \end{pmatrix}\\
    \exists Q \text{ inversible tq } A &= Q \begin{pmatrix}
        0 & & &  \\
          & \ddots & & & \\
          & & 0 & & \\
          & & & \lambda_1 & \\
          & & & & \lambda_2 & \\
          & & & & & & \ddots
    \end{pmatrix}Q^{-1}\\
    \text{Donc } A^2 &= Q \begin{pmatrix}
        0 & & &  \\
          & \ddots & & & \\
          & & 0 & & \\
          & & & \lambda_1^2 & \\
          & & & & \lambda_2^2 & \\
          & & & & & & \ddots
    \end{pmatrix}Q^{-1}
\end{align*}
\begin{align*}
    \lambda_1 + \lambda_2 &= trA = 2 \\
    \lambda_1^2 + \lambda_2^2 &= 2(n-2) + 2n \\
    &= 4(n-1)\\
    (\lambda_1 + \lambda_2)^2 &= \lambda_1^2 + \lambda_2^2 + 2\lambda_1 \lambda_2\\
    \lambda_1 \lambda_2 &= \frac{1}{2}((\lambda_1 + \lambda_2)^2 - (\lambda_1^2 + \lambda_2^2))\\
    &\begin{cases}
        \lambda_1 + \lambda_2 = 2\\
        \lambda_1 \lambda_2 = 4-2n
    \end{cases}\\
    &\Leftrightarrow \lambda_1 \lambda_2 \text{ sont les racines de } X^2 - 2X + 4 - 2n\\
    \Delta &= 4 - 4(4-2n)\\
    &= 4(2n-3)>0\\
  \Rightarrow \{\lambda_1, \lambda_2\} &= \{1 + \sqrt{2n-3}, 1-\sqrt{2n-3}\}
\end{align*}
a4/ \(Sp(A) = \{\lambda_1, \lambda_2, 0\}\)\\
\begin{equation*}
    (X-\lambda_1)(X-\lambda_2)X \big| \Pi_A \big| \chi_A = (X-\lambda_1)(X-\lambda_2)X^{n-2}
\end{equation*}
A diagonalisable\\
Donc \(\Pi_A\) scindé simple\\
Donc \(\Pi_A = X^3 - 2X^2 + 4(-2n)X\)\\
\(\Pi_A = 0\)\\
Donc \(A^3 - 2A^2 + (4-2n)A = 0\)\\
On suppose \(A^k = \lambda_k A + \mu_k A^2\)\\
\(A^{k+1} = \lambda_k A^2 + \mu_k A^3 \in Vect(A, A^2)\) \\

\noindent 
\(A^k ?\)
\(\Pi_A (A) = 0\)\\
division euclidienne de \(X^k\) par \(\Pi_A\)\\
\(\begin{cases}
    X^k = Q\Pi_A + R \\
    R \in \mathbb{R}_2[X]
\end{cases}\)\\
je vois pas le tableau\\
Donc \begin{equation*}
    A^k = (Q\Pi_A)(A) + \lambda_k A^2 + \mu_k A + \nu_k I_n
\end{equation*}\\
Donc \begin{align*}
    X^k = Q\Pi_A + &\lambda_k X^2 + \mu_k X + \nu_k \\ 
    \text{On évalue en 0 : }\\
    0 &= Q\Pi_A(0) + \nu_k\\
    &\Leftrightarrow \nu_k = 0
\end{align*}


\noindent \colorbox{green}{260.5}\\
\underline{\(\varphi\) est linéaire} \\
par linéarité de la dérivée et la bilinéarité du produit \\
\underline{\(\varphi\) est à valeurs dans \(\mathbb{R}_4[X]\)}\\
\begin{align*}
    \varphi(X^k) &= (X^2 -1)X^{k-1} - (4X+1)X^k\\
    &= (k-4)X^{k+1} - X^k - kX^{k-1}\\
    &\in \mathbb{R}_4[X]\\
    \varphi(1) &\in \mathbb{R}_4[X]
\end{align*}
b/ Soit \begin{align*}
    I_1 &= ]-\infty, -1[\\
    I_2 &= ]-1, 1[\\
    I_3 &= ] 1, +\infty[
\end{align*}
Sur l'intervalle \(I_k\)\\
C'est une EQD linéaire d'ordre 1 homogène \\
\begin{align*}
    S = \begin{cases}
        I_r \to \mathbb{R}\\
        x \mapsto A|x-1|^{\frac{5+\lambda}{2}}|x+1|^{\frac{3-\lambda}{2}}
    \end{cases}
\end{align*}
c/ \begin{align*}
    \varphi(P) &= \lambda P\\
    &\Leftrightarrow (X^2 -1)P'-(4X+1)P = \lambda P \\
    &\Leftrightarrow \forall x \in ]1, +\infty[\\
    &(x^2-1)P'(x) - (4x+1)P(x) = \lambda P(x)\\
    &\Leftrightarrow \forall x \in ]1, +\infty[ \\
    P'(x) &= \frac{4x+1+\lambda}{x^2-1}P(x)\\
    &= (\frac{5+\lambda}{2(x-1)} + \frac{3-\lambda}{2(x+1)}) P(x)\\
    &\Leftrightarrow P \text{ sol de } E_\lambda \text{ de la question b }\\
    &\Leftrightarrow \exists A \text{ tq } P = A(X-1)^{\frac{5+\lambda}{2}}(X+1)^{\frac{3-\lambda}{2}}
\end{align*}\\
En posant \(\lambda\) tq \begin{align*}
    \frac{5+\lambda}{2}, \frac{3-\lambda}{2} \in \mathbb{N}\\
    \frac{5+\lambda}{2} + \frac{3-\lambda}{2} \leq 4
\end{align*}
\begin{align*}
    \lambda = -5 \\
    \lambda = -3 \\
    \lambda = -1 \\
    \lambda = 1 \\
    \lambda = 3 
\end{align*}

\noindent \colorbox{green}{260.27}\\
\begin{align*}
    A &= \begin{pmatrix}
        0 & 1 & 0 \\ 1& 0 & 1\\ 2& 0 & 0
    \end{pmatrix} &B = \begin{pmatrix}
        0 & 1 & 1\\ 1& 0 & 0\\ 0 & 2 & 0
    \end{pmatrix}\\
    \chi_A &= \begin{vmatrix}
        X&-1&0\\-1&X&-1\\-2&0&X
    \end{vmatrix} &\chi_B = \begin{vmatrix}
        X&-1&-1\\-1&X&0\\0&-2&X
    \end{vmatrix}\\
    \chi_A &= X^3 - X - 2 &\chi_B = X^3 - X - 2
\end{align*}
b/ \begin{align*}
    f:x &\mapsto x^3 - x - 2\\
    f':x &= 3x^2 -1\\
    &= 3(x-\frac{1}{\sqrt{3}})(1+\frac{1}{\sqrt{3}})\\
    f(-\frac{1}{\sqrt{3}}) &= \frac{1}{3^\frac{3}{2}} - \frac{1}{\sqrt{3}}-2 < 0
\end{align*}
Veuillez imaginer un beau tableau de variations j'ai la flemme\\
f admet une unique racine \(\alpha\) \begin{align*}
    \chi_A = (X-\alpha)(X-\beta)(X-\overline{\beta}) \indent \beta \notin \mathbb{R}
\end{align*}
non scindé dans \(\mathbb{R}\) \\
Alors A et B ne sont pas diagonalisables\\
\(\chi_A\) est scindé simple dans \(\mathbb{R}\)\\
A et B sont semblables dans \(\mathcal{M}_3(\mathbb{C})\)
à \(\begin{pmatrix}
    \alpha & & \\ & \beta & \\ & & \overline{\beta}
\end{pmatrix}\) donc sont semblables entre elles dans \(\mathbb{C}\)\\

\noindent \colorbox{green}{260.30}\\
Montrons que : \begin{equation*}
    \chi_A(B) \in GL_n(\mathbb{K}) \Leftrightarrow Sp_\mathbb{C}(A)\cap Sp_\mathbb{C}(B) = \oslash 
\end{equation*}\\
\fbox{\(\Rightarrow\)} Contraposée : supposons \(Sp_\mathbb{C}(A)\cap Sp_\mathbb{C}(B) \neq \oslash \)
Donc \(\exists \lambda \in \mathbb{C}\) vp de A et de B : \begin{align*}
    \chi_A &= (X-\lambda)Q_A\\
    \chi_B &= (X-\lambda)Q_B\\
    \chi_A(B) &= (B-\lambda I_n)Q_A(B)\\
    det(\chi_A(B)) &= det(B-\lambda I_n)det(Q_A(B))\\
    &= 0 \text{ car \(\lambda\) vp de B}
\end{align*}
\fbox{\(\Leftarrow\)} Contraposée : \(det(\chi_A(B)) = 0\) \\
\(\chi_A \in \mathbb{C}[X]\)\\
Donc \(\chi_A\) scindé : \(\chi_A = \prod_{k=1}^{p}(X-\lambda_k)^{m(\lambda_k)}\), 
\((\lambda_k)_k\) exactement les vp de A\\
Donc \begin{align*}
    0 &= det(\chi_A(B))\\
    &= det(\prod_{k=1}^{p}(B-\lambda_k)^{m(\lambda_k)})\\
    &= \prod_{k=1}^{p}det((X-\lambda_k)^{m(\lambda_k)})
\end{align*}
Donc \(\exists k \in [|1,p|]\) tq \(det(B-\lambda_k I_n) = 0\)
donc \(\lambda_k\) vp de B \\
Alors \(Sp_\mathbb{C}(A)\cap Sp_\mathbb{C}(B) \neq \oslash \)

\noindent \colorbox{green}{260.21}\\
\begin{align*}
    A = \begin{pmatrix}
        -1 & 1 & -3\\ 0 & 2 & -1 \\ 1& 0 & 2
    \end{pmatrix}\\
    \chi_A = (X-1)^3
\end{align*}
Si A est diagonalisable alors \(A = PIdP^{-1} = Id\) absurde\\
Donc A n'est pas diagonalisable\\
b/ On a \(B = A -I_3\)\\
Par le théorème de Cayley-Hamilton \(\chi_A(A) = 0\) donc \(B^3 = 0\)\\
\begin{align*}
    B &= \begin{pmatrix}
        -2 & 1 & -3\\0& 1& -1\\1&0&1
    \end{pmatrix} \text{ donc } &B^2 = \begin{pmatrix}
        1&-1&2\\-1&1&-2\\-1&1&-2
    \end{pmatrix}
\end{align*}
Donc \(\exists X \in \mathcal{M}_{3,1}(\mathbb{R})\) tq \(X \notin KerB\)\\
Soit \(\lambda_1, \lambda_2, \lambda_3\) tq \(\lambda_1 X + \lambda_2 B X + \lambda_3 B^2 X = 0\)\\
\(\lambda_1B^2 X = 0\) donc \(\lambda_1 = 0\)\\
En multipliant par B on a \(\lambda_2\) = 0\\
Donc \(\lambda_3 B^2 X = 0\)\\
Donc  \(\lambda_3 = 0\)\\
Ainsi \((X, BX, B^2X)\) est une base de \(\mathcal{M}_{3,1}(\mathbb{R})\)\\
\begin{align*}
    Mat(b, \mathcal{B}) &= \begin{pmatrix}
        0 & 1 & 0 \\ 0 & 0 & 1\\ 0 & 0 & 0 
    \end{pmatrix}\\
    P&= Pass(\mathcal{B}_\mathcal{C} \to \mathcal{B})
\end{align*}
On a \(\begin{pmatrix}
    1\\0\\0
\end{pmatrix} \notin KerB\)\begin{align*}
    Pass(\mathcal{B}_\mathcal{C} \to \mathcal{B}) &= \begin{pmatrix}
        1 & -2 & 1\\ -1 & 0 & 0 \\ -1 & 1 & 0
    \end{pmatrix}\\
    B &= PB'P^{-1}\\
    A-I_3 &= PB'P{-1}\\
    A &= PB'P{-1} + I_3\\
    &= P(B'+I_3)P^{-1} \indent \text{ car } I_3 = P I_3 P^{-1}\\
    B'+I_3 &= \begin{pmatrix}
        1 & 1& 0\\ 0 & 1& 1\\ 0& 0 & 1
    \end{pmatrix}
\end{align*}




\end{document}
